\documentclass{ximera}
\title{MATH 3307: Linear Algebra\\
Project 3: Kaleidoscope Extended Project}
\author{Ahniya Muse}
\date{April 2026}

\begin{document}

\maketitle

\section*{Overview}

The goal of this project is to use matrices and linear transformations to model a virtual kaleidoscope. The transformations used are reflections and rotations in the plane.


\centering
\includegraphics[width=0.28\textwidth]{start.jpg}
\includegraphics[width=0.28\textwidth]{horizontal-reflection.jpg}
\includegraphics[width=0.28\textwidth]{diagonal-reflection.jpg}

\vspace{0.3cm}

\includegraphics[width=0.45\textwidth]{finish.jpg}
\caption{Starting image, reflected images, and completed kaleidoscope image.}
\end{figure}

\section*{(1) Reflection Across the x-axis}

The matrix for reflection across the x-axis is:

\[
F =
\begin{bmatrix}
1 & 0 \\
0 & -1
\end{bmatrix}
\]

This transformation keeps the x-coordinate the same and changes the sign of the y-coordinate.

\section*{(2) Rotation Counterclockwise by $60^\circ$}

The matrix for rotation by $60^\circ$ is:

\[
R =
\begin{bmatrix}
\frac{1}{2} & -\frac{\sqrt{3}}{2} \\
\frac{\sqrt{3}}{2} & \frac{1}{2}
\end{bmatrix}
\]

This rotates every point $60^\circ$ counterclockwise about the origin.

\section*{(3) Reflection Across the Line at $60^\circ$}

The matrix for reflection across the line at $60^\circ$ is:

\[
RF =
\begin{bmatrix}
-\frac{1}{2} & \frac{\sqrt{3}}{2} \\
\frac{\sqrt{3}}{2} & \frac{1}{2}
\end{bmatrix}
\]

This matrix is obtained by reflecting across the x-axis and then rotating by $60^\circ$.

\section*{(4a) Composition and Non-Commutativity}

\[
RF =
\begin{bmatrix}
-\frac{1}{2} & \frac{\sqrt{3}}{2} \\
\frac{\sqrt{3}}{2} & \frac{1}{2}
\end{bmatrix}
\]

\[
FR =
\begin{bmatrix}
\frac{1}{2} & -\frac{\sqrt{3}}{2} \\
-\frac{\sqrt{3}}{2} & -\frac{1}{2}
\end{bmatrix}
\]

Since $RF \neq FR$, matrix multiplication is not commutative.

This means the order of transformations matters. Reflecting first and then rotating gives a different result than rotating first and then reflecting.

\section*{(4b) Powers of Transformations}

\[
R^2 =
\begin{bmatrix}
-\frac{1}{2} & -\frac{\sqrt{3}}{2} \\
\frac{\sqrt{3}}{2} & -\frac{1}{2}
\end{bmatrix}
\]

\[
R^3 =
\begin{bmatrix}
-1 & 0 \\
0 & -1
\end{bmatrix}
\]

\[
R^6 =
\begin{bmatrix}
1 & 0 \\
0 & 1
\end{bmatrix}
\]

The matrix $R^6 = I$ because rotating six times by $60^\circ$ results in a full $360^\circ$ rotation, returning the image to its original position.

\section*{(4c) Determinants as Geometric Invariants}

\[
\det(R) = 1
\]

Rotation preserves both area and orientation.

\[
\det(F) = -1
\]

Reflection preserves area but reverses orientation.

\[
\det(RF) = -1
\]

This preserves area but reverses orientation.

\vspace{0.3cm}

All kaleidoscope transformations preserve distances because rotations and reflections are rigid motions, meaning they do not change the size or shape of the image.

\end{document}